\subsection{Grasper Verification}

\subsubsection{Holding Capacity}

The results in table \ref{tab:weight-inc} show that the device is capable of holding
over twice the target weight. Using equation \ref{eqn:load-lim}, the verified 
limit of 8.7 kg means this device as currently designed can support a drone of up to 1.17 kg.


\subsubsection{Weight}

The total weight of the device is well over the target of 47.5g; the achieved weight of
236g is almost 5x the target. There are a number of opportunities for improvement on this
however: 

\begin{itemize}
    \item The results in Section \ref{sec:results:capacity} show that the holding capacity 
        is also significantly over the target, so the device could be used as-is on a larger drone with 
        more payload capacity.
    \item Based on the results in Section \ref{sec:results:capacity}, the device could also be made
        smaller and lighter, though weaker, to only support a 500g drone.
    \item The device currently is designed as an isolated component. If it were integrated more tightly
        into a drone's airframe, the weight of the baseplate could be eliminated.
    \item Similarly, the electronics could be more tightly integrated to reduce the number and size of
        PCBs and connectors used. The Raspberry Pi used for image processing may also perform some flight
        control functions, reducing the electronics weight needed for the drone.
\end{itemize}


\subsubsection{Repeatability}

The repeatability tests reveal some issues with the grasper control logic, specifically 
with detecting when the grasper has reached its limits based on the motor current. The first 
issue, premature triggering of the current limit, seems simple to solve by raising the threshold
current. However, when releasing the grasper it was observed that there is typically a 10-15 second
period after the grasper reaches its fully released position during which the current limit does not
trigger. Raising the current threshold could extend this period further, increasing wear on the device,
or prevent it from detecting the open position entirely. 

The second issue, failure to detect the released position, seems to indicate an issue with drifting 
servo tuning. That the failure did not occur until the end of the test, all instances occurred 
sequentially, and the failure did not occur again after re-tuning the zero point, all suggest that
the issue is not due to random errors in the current detection, but rather a systematic flaw in 
the controller. It may be that repeatedly driving the servo while stalled accelerates drift in the 
zero point. This could be solved by driving the device directly with DC motors rather than servos,
which would eliminate the potential for zero point drift entirely. This could also provide opportunities
for weight reduction. 

Overall the results in Section \ref{sec:results:repeatability} show that thought the grasper concept is functional, 
it needs work to improve reliability. Swapping the servos for DC motors may provide some improvement, 
and eliminating the current based limit detection in favor of limit switches or optical 
sensors to detect the limits of travel may also be needed.

